\documentclass{article}
\usepackage[utf8]{inputenc}
\usepackage[T1]{fontenc}
\usepackage[margin=2cm]{geometry}
\usepackage{hyperref}
\sloppy
\begin{document}
\section*{Historical source}
Encyclopaedia Britannica scan, source row 560. Parse 3.5 extracted the text. Super 3.5 VL corrected it from an enlarged crop of each block. It can still contain recognition errors.
Page headers, footers, and page numbers are removed.

\subsection*{Untitled continuation (E0047)}
\textit{Continues an entry from an earlier page.}

Amethy- gular, semiquinquefid, and perfluent: The corolla is monopetalous; the border quinquepartite, the lowest division more expanding: The filamenta consist of two slender filaments approximated; the anthers are simple and roundish: The pifillium has a four-cleft germen; stylus, the size of the flamina; stigmata two, acute: No corolla: The feeds four, gibbous, and shorter than the calyx:—There is only one known species.

This plant is a native of Siberia, from whence the feeds were sent to the imperial garden at Petersburg, and thence brought to Britain. It is an annual plant, with an upright stalk, which rises about a foot high. Towards the top it puts forth two or three small lateral branches, garnished with small trifid leaves, sawed on their edges, of a very dark green colour. The flowers appear in June or July, and are produced in small umbels at the extremities of the branches. They are of a fine blue colour, as are also the upper part of the branches, and the leaves immediately under the umbel, so that they make a fine appearance.

Culture. The feeds of the amethystea should be sown in autumn, as they are apt to remain a whole year in the ground if kept till the spring. When the plants come up, nothing else is necessary than to keep them clear of weeds, and to thin them where they are too close. They ought to be sown where they are to remain, as they do not thrive when transplanted.

AMETHYSTINE is applied, in antiquity, to a kind of purple garment dyed of the hue of amethyst. In this sense amethystine differed from Tyrian as well as from hyacinthine purple, being a kind of medium between both.

\subsection*{AMHAR (E0048)}
AMHAR, or AMHARA, a province of Abyssinia, said to extend 40 leagues from east to west. It is considered as the most noble in the whole empire, both on account of its being the usual residence of the Abyssinian monarchs, and having a particular dialect different from all the rest, which, by reason of the emperors being brought up in this province, is become the language of the court and of the politer people. Here is the famed rock Amba-geihen, where the young monarchs were formerly confined. See Amba.

AMHURST (Nicholas), an English poet and political writer, was born at Marden in Kent, and entered of St. John's college, Oxford; from whence he was expelled for irregularity of conduct and libertine principles. Retaining great resentment against the university on this account, he abused its learning and discipline, and some of the most respectable characters in it, in a poem published in 1724 called Oculus Britanniae, and in a book entitled Terra Filius. He published, A Miscellany of Poems, sacred and profane; and, The Convocation, a poem in five cantos, which was a satire on the Bishop of Bangor's antagonists. But he is best known for the share he had in the political paper called The Craftsman; though, after having been the drudge of his party for near 20 years, he was as much forgot in the famous compromise of 1742 as if he had never been born; and, when he died in that year of a broken heart, was indebted to the charity of his bookseller for a grave.

a grave.
AMIANTHUS, or EARTH-FLAX, in natural hi-
story, a fibrous, flexible, elastic, mineral substance, con-
sisting of short, abrupt, and interwoven filaments. It
is found in Germany, in the strata of iron ore, some-

Amianthus is called by Mr Kirwan in the muriatic genus of earths, because it contains about a fifth part of magnesia. Its other constituents are, flint, mild calcareous earth, barytes, clay, and a very small proportion of iron. It is fusible per fe in a strong heat, and also with the common fluxes. It differs from asbestos in containing some ponderous earth.

\subsection*{AMICABLE (E0049)}
AMICABLE, in a general sense, denotes any thing done in a friendly manner, or to promote peace.

Amicable Benches, in Roman antiquity, were, according to Pitiscus, lower and less honourable feats allotted for the judices pedanei, or inferior judges, who, upon being admitted of the emperor's council, were dignified by him with the title amici.

Amicable Numbers, such as are mutually equal to the sum of one another's aliquot parts. Thus the numbers 284 and 220 are amicable numbers: for the aliquot parts 1, 2, 4, 5, 10, 11, 20, 22, 44, 55, 110, of 220, are together equal to the other number 284; and the aliquot parts 1, 2, 4, 71, 142, of 284, are together equal to 220.

\subsection*{AMICTUS (E0050)}
AMICTUS, in Roman antiquity, was any upper garment worn over the tunica.

Amictus, among ecclesiastical writers, the uppermost garment anciently worn by the clergy : the other five being the alba, fingulum, sola, manipulus, and planeta. The amictus was a linen garment, of a square figure, covering the head, neck, and shoulders, and buckled or clasped before the breast. It is still worn by the religious abroad.

by the religious abroad.
AMICULUM, in Roman antiquity, a woman's up-
per garment, which differed from the pala. It was
worn both by matrons and courtezans.
AMICUS curiae, a law term, to denote a by-stand-

AMICUS curiæ, a law term, to denote a by-stander who informs the court of a matter in law that is doubtful or mistaken.

\subsection*{AMIDA (E0051)}
AMIDA, a god worshipped by the Japanese, who has many temples erected to him in the island of Japan, of which the principal is at Jedo. The Japanese have such a confidence in their idol Amida, that they hope to attain eternal felicity by the frequent invocation of his name. One of the figures of this idol is represented at Rome.

Amida (anc. geog.), a principal city of Mesopotamia (Liber Notitiae); Ammaea (Ptolemy); situated on a high mountain, on the borders of Afyria, on the Tigris, where it receives the Nymphius.—It was taken from the Romans, in the time of the emperor Constantius, by Saporas king of Persia. The siege is said to have cost him 20,000 men; however, he reduced it to such ruin, that the emperor afterwards wept over it. According to Ammianus Marcellinus, the city was raided; the chief officers were crucified; and the rest, with the soldiers and inhabitants, either put to the sword or carried into captivity, except our historian himself, and two or three more, who, in the dead of the night, escaped through a postern unperceived by the enemy. The inhabitants of Nisibis, however, being obliged to leave their own city by Jovian's treaty with


\section*{Source record}
\begin{description}
\item[Scan ID] \texttt{britannica:309f765b1ecc83e0a7f7762e44b4a927a7555443:train:560}
\item[Image file] \texttt{page-0560-c458f594fc79.jpg}
\item[Image SHA-256] \texttt{c458f594fc79392e42082ac7a8c338082e8d5e2c2c79dd2f49f1940b91144300}
\item[Dataset] NationalLibraryOfScotland/encyclopaedia\_britannica\_illustrated at revision \texttt{309f765b1ecc83e0a7f7762e44b4a927a7555443} (CC0-1.0)
\item[Extraction] model \texttt{nvidia/nemotron-3.5-parse}, raw reply \texttt{extract/raw/page-0560-c458f594fc79.json}
\item[Raw reply SHA-256] \texttt{c3ff53713741c121b2aad52d415d12aefbbbbe63366d78fce517442f41fa6082}
\item[Correction] model \texttt{nvidia/Nemotron-3.5-Super-VL}, block crops, raw replies \texttt{correct/raw/crop\_nothink/page-0560-c458f594fc79-b*.json}; super\_vl: 17
\item[Running header] AMI [ 619] A M I
\end{description}
\end{document}
